\section{Second-moment tensor and the dihedral case}
  \label{sec:M}

  \begin{definition}
    The \emph{second-moment tensor} of the generator set $S$ is
    \begin{equation}
      M \;:=\; \sum_{s \in S^+} \vec{u}_s\,\vec{u}_s^{\,T} \;\in\; \mathrm{Sym}_3(\R).
      \label{eq:M-def}
    \end{equation}
  \end{definition}

  By construction $M$ is positive semidefinite.
  Its isotropy class --- whether $M \propto I$ or not --- is the first
  character-theoretic discriminant in the classification.

  \begin{proposition}[Dihedral anisotropy]
    \label{prop:dic-anisotropy}
    Let $\rho(G) \cong \mathrm{Dic}_n$ with $n \ge 3$, and let $S$ be the full set of
    traceless elements of $G$.
    In the basis where $\vec{e}_z$ is the $C_{2n}$ symmetry axis, the second-moment
    tensor $M$ has eigenvalues
    \begin{equation}
      \mathrm{spec}(M) \;=\;
      \begin{cases}
        (n,\; n,\; 0)   & \text{if $n$ is odd,} \\
      (n,\; n+2,\; 0) & \text{if $n$ is even.}
      \end{cases}
      \label{eq:dic-eigenvalues}
    \end{equation}
    In particular, $M \not\propto I$ for all $n \ge 3$.
  \end{proposition}

  \begin{proof}
    In the standard presentation, $\mathrm{Dic}_n$ has a cyclic normal subgroup
    $\langle a \rangle \cong \Z_{2n}$ generated by
    $a = (\cos(\pi/n),\, \sin(\pi/n),\, 0,\, 0)$
    and an order-4 element $b = (0,\, 0,\, 1,\, 0)$ satisfying $bab^{-1} = a^{-1}$.
    We identify quaternions with $(a_0, a_1, a_2, a_3)$ and the corresponding axis
    via $(0, a_1, a_2, a_3) \mapsto \vec{u} = (a_2, a_3, a_1)$.

    \medskip
    \textbf{(A) Flip-coset elements.}
    A direct computation gives
    \begin{equation}
      b\,a^k \;=\; (0,\; \sin(k\pi/n),\; \cos(k\pi/n),\; 0),
      \qquad k = 0, \ldots, 2n-1,
      \label{eq:dic-flip}
    \end{equation}
    hence unit axes $\vec{u}_k = (\cos(k\pi/n),\, 0,\, \sin(k\pi/n))$.
    All $2n$ flip-coset axes are uniformly distributed in the $xz$-plane.
    Using the standard identities for $2n$ uniform points on a circle,
    \[
      \sum_{k=0}^{2n-1} \cos^2\!\tfrac{k\pi}{n} = n, \quad
      \sum_{k=0}^{2n-1} \sin^2\!\tfrac{k\pi}{n} = n, \quad
      \sum_{k=0}^{2n-1} \cos\tfrac{k\pi}{n}\sin\tfrac{k\pi}{n} = 0,
    \]
    the flip-coset contribution to $M$ is
    \[
      M_{\mathrm{flip}}
      \;=\; \sum_{k=0}^{2n-1} \vec{u}_k\,\vec{u}_k^T
      \;=\; \diag(n,\; 0,\; n).
    \]

    \medskip
    \textbf{(B) Cyclic traceless elements ($n$ even only).}
    The cyclic element $a^k$ has real part $\cos(k\pi/n)$, which vanishes exactly for
    $k = n/2$ and $k = 3n/2$ (when $n$ is even, and not otherwise).
    These two elements are $a^{n/2} = (0, 1, 0, 0) = \pm i$ with axis $\vec{e}_z = (0,0,1)$,
    contributing $2\,\vec{e}_z\vec{e}_z^T = \diag(0, 0, 2)$.

    \medskip
    \textbf{(C) Total.}
    Combining both contributions:
    \[
      M \;=\;
      \begin{cases}
        \diag(n,\; 0,\; n)   & \text{$n$ odd (no cyclic traceless elements),} \\
        \diag(n,\; 0,\; n+2) & \text{$n$ even (two cyclic traceless elements $\pm i$).}
      \end{cases}
    \]
    In both cases the eigenvalues are $(n,\, 0,\, n)$ or $(n,\, 0,\, n+2)$,
    i.e.\ the spectrum stated in~\eqref{eq:dic-eigenvalues} (reordering the diagonal).
    Since the zero eigenvalue is strictly less than $n \ge 3$, we have $M \not\propto I$.
  \end{proof}

  \begin{remark}
    \label{rem:Dic2-iso}
    The special case $n = 2$ (i.e.\ $\mathrm{Dic}_2 \cong Q_8$) is \emph{isotropic}:
    the six traceless elements $\{\pm i, \pm j, \pm k\}$ span all three coordinate axes,
    giving $M_{Q_8} = 2I$.
    This is not contradicted by~\eqref{eq:dic-eigenvalues}: the formula applies to
    $\mathrm{Dic}_n$ for $n \ge 2$, but for $Q_8 = \mathrm{Dic}_2$ the cyclic subgroup
    contributes traceless elements $\pm i$ (axis $\vec{e}_z$) \emph{and} the flip coset
    contributes $\pm j, \pm k$ (axes $\vec{e}_x, \vec{e}_y$), restoring full isotropy.
    Proposition~\ref{prop:dic-anisotropy} is thus sharp: the anisotropy holds exactly for
    $n \ge 3$.
  \end{remark}
